\chapter{2005年湖北卷（文）}
\section{单选题}
\begin{problem}
设 \(P\)，\(Q\) 为两个非空实数集合，定义集合 \(P + Q = \{a + b \mid a \in P, b \in Q\}\)，若 \(P = \{0, 2, 5\}\)，\(Q = \{1, 2, 6\}\)，则 \(P + Q\) 中元素的个数是（\quad）

\choices
    {\(9\)}
    {\(8\)}
    {\(7\)}
    {\(6\)}
\end{problem}

\begin{answer}
B
\end{answer}

\begin{solution}
由集合加法的定义，\(P+Q\) 中的和为
\(0+1=1\)，\(0+2=2\)，\(0+6=6\)，\(2+1=3\)，\(2+2=4\)
\(2+6=8\)，\(5+1=6\)，\(5+2=7\)，\(5+6=11\).
去掉重复的 \(6\)，得 \(P+Q=\{1,2,3,4,6,7,8,11\}\)，共有 \(8\) 个元素，故选 B.
\end{solution}

\begin{problem}
对任意实数 \(a\)，\(b\)，\(c\) 给出下列命题：

\begin{circlelist}
\item “\(a = b\)”是“\(ac = bc\)”充要条件；
\item “\(a + 5\) 是无理数”是“\(a\) 是无理数”的充要条件；
\item “\(a > b\)”是“\(a^{2} > b^{2}\)”的充分条件；
\item “\(a < 5\)”是“\(a < 3\)”的必要条件.
\end{circlelist}

其中真命题的个数是（\quad）

\choices
    {\(1\)}
    {\(2\)}
    {\(3\)}
    {\(4\)}
\end{problem}

\begin{answer}
B
\end{answer}

\begin{solution}
第1个命题中，\(a=b\) 能推出 \(ac=bc\)，但取 \(c=0\) 且 \(a\ne b\) 时，逆命题不成立，故第1个命题是假命题. 由于 \(5\) 是有理数，\(a+5\) 是无理数当且仅当 \(a\) 是无理数，所以第2个命题是真命题. 第3个命题不成立，例如 \(a=1\)，\(b=-2\) 时，\(a>b\)，但 \(a^2<b^2\). 第4个命题是真命题，因为 \(a<3\) 必能推出 \(a<5\)，即 \(a<5\) 是 \(a<3\) 的必要条件. 因此真命题有 \(2\) 个，故选 B.
\end{solution}

\begin{problem}
已知向量 \(\bs{a} = (-2, 2)\)，\(\bs{b} = (5, k)\). 若 \(\left|\bs{a} + \bs{b}\right|\) 不超过 5，则 \(k\) 的取值范围是（\quad）

\choices
    {\([-4, 6]\)}
    {\([-6, 4]\)}
    {\([-6, 2]\)}
    {\([-2,6]\)}
\end{problem}

\begin{answer}
C
\end{answer}

\begin{solution}
由 \(\bs{a}+\bs{b}=(3,k+2)\)，得
\[
\left|\bs{a}+\bs{b}\right|\leqslant5\Longleftrightarrow 3^2+(k+2)^2\leqslant25\Longleftrightarrow (k+2)^2\leqslant16
\]
所以 \(-4\leqslant k+2\leqslant4\)，即 \(-6\leqslant k\leqslant2\)，故选 C.
\end{solution}

\begin{problem}
函数 \( y = \e^{|\ln x|} - |x - 1| \) 的图象大致是（\quad）

\choices
    {\choicebitmap[width=0.15\paperwidth]{img/2005/hubei_liberal/q4_optA_clean.png}}
    {\choicebitmap[width=0.15\paperwidth]{img/2005/hubei_liberal/q4_optB_clean.png}}
    {\choicebitmap[width=0.15\paperwidth]{img/2005/hubei_liberal/q4_optC_clean.png}}
    {\choicebitmap[width=0.15\paperwidth]{img/2005/hubei_liberal/q4_optD_clean.png}}
\end{problem}

\begin{answer}
D
\end{answer}

\begin{solution}
函数的定义域为 \(x>0\). 当 \(0<x<1\) 时，\(|\ln x|=-\ln x\)，\(|x-1|=1-x\)，从而
\[
y=\e^{-\ln x}-(1-x)=\frac1x+x-1
\]
此时 \(y'=1-\frac1{x^2}<0\)，且当 \(x\to0^+\) 时 \(y\to+\infty\)，当 \(x\to1^-\) 时 \(y\to1\). 当 \(x\geqslant1\) 时，\(|\ln x|=\ln x\)，\(|x-1|=x-1\)，所以 \(y=x-(x-1)=1\). 因此图象在 \((0,1)\) 上从无穷大递减到 \((1,1)\)，在 \([1,+\infty)\) 上为直线 \(y=1\)，故选 D.
\end{solution}

\begin{problem}
木星的体积约是地球体积的 \(240 \sqrt{30}\) 倍，则它的表面积约是地球表面积的（\quad）

\choices
    {\(60\) 倍}
    {\(60 \sqrt{30}\) 倍}
    {\(120\) 倍}
    {\(120 \sqrt{30}\) 倍}
\end{problem}

\begin{answer}
C
\end{answer}

\begin{solution}
设木星与地球的半径之比为 \(k\)，则体积之比为 \(k^3\)，表面积之比为 \(k^2\). 由题意，
\[
k^3=240\sqrt{30}=(2\sqrt{30})^3=(\sqrt{120})^3
\]
所以 \(k=\sqrt{120}\)，从而表面积之比为 \(k^2=120\)，故选 C.
\end{solution}

\begin{problem}
双曲线 \(\frac{x^2}{m} - \frac{y^2}{n} = 1\)（\(mn \neq 0\)）离心率为 \(2\)，有一个焦点与抛物线 \(y^2 = 4x\) 的焦点重合，则 \(mn\) 的值为（\quad）

\choices
    {\(\frac{3}{16}\)}
    {\( \frac{3}{8} \)}
    {\(\frac{16}{3}\)}
    {\(\frac{8}{3}\)}
\end{problem}

\begin{answer}
A
\end{answer}

\begin{solution}
抛物线 \(y^2=4x\) 的焦点为 \((1,0)\). 由于双曲线有焦点与之重合，双曲线的焦点在 \(x\) 轴上，故可设其标准形式中的半长轴、半短轴和半焦距分别为 \(a\)，\(b\)，\(c\)，且 \(m=a^2>0\)，\(n=b^2>0\). 由焦点重合得 \(c=1\)，由离心率得
\[
\frac{c}{a}=2\Longrightarrow a=\frac12\Longrightarrow m=a^2=\frac14
\]
又因为 \(c^2=a^2+b^2\)，所以
\[
n=b^2=c^2-a^2=1-\frac14=\frac34
\]
因此 \(mn=\frac14\cdot\frac34=\frac{3}{16}\)，故选 A.
\end{solution}

\begin{problem}
在 \( y = 2^{x} \)，\( y = \log_{2} x \)，\( y = x^{2} \)，\( y = \cos 2x \) 这四个函数中，当 \( 0 < x_{1} < x_{2} < 1 \) 时，使 \( f\left(\frac{x_{1} + x_{2}}{2}\right) > \frac{f(x_{1}) + f(x_{2})}{2} \) 恒成立的函数的个数是（\quad）

\choices
    {\(0\)}
    {\(1\)}
    {\(2\)}
    {\(3\)}
\end{problem}

\begin{answer}
B
\end{answer}

\begin{solution}
对四个函数分别考察其在 \((0,1)\) 上的凹凸性：\(2^x\) 与 \(x^2\) 的二阶导数分别为正，均为凸函数，故中点处函数值小于两端函数值的平均数；\(\log_2x\) 的二阶导数为
\[
\left(\log_2x\right)''=-\frac{1}{x^2\ln2}<0
\]
所以它在该区间上为严格凹函数，题中的不等式恒成立. 对于 \(\cos2x\)，有 \(\left(\cos2x\right)''=-4\cos2x\). 当 \(\frac\pi4<x<1\) 时，\(\cos2x<0\)，故 \(\left(\cos2x\right)''>0\)，函数在该区间上严格凸；取 \(\frac\pi4<x_1<x_2<1\)，则中点处函数值小于两端函数值的平均数，题中的不等式不成立. 因此只有 \(\log_2x\) 满足条件，函数个数为 \(1\)，故选 B.
\end{solution}

\begin{problem}
已知 \(a\)，\(b\)，\(c\) 是直线，\(\beta\) 是平面，给出下列命题：

\begin{circlelist}
\item 若 \(a \perp b\)，\(b \perp c\)，则 \(a \parallel c\)；
\item 若 \(a \parallel b\)，\(b \perp c\)，则 \(a \perp c\)；
\item 若 \(a \parallel \beta\)，\(b \subset \beta\)，则 \(a \parallel b\)；
\item 若 \(a\) 与 \(b\) 异面，且 \(a \parallel \beta\)，则 \(b\) 与 \(\beta\) 相交；
\item 若 \(a\) 与 \(b\) 异面，则至多有一条直线与 \(a\)，\(b\) 都垂直.
\end{circlelist}

其中真命题的个数是（\quad）

\choices
    {\(1\)}
    {\(2\)}
    {\(3\)}
    {\(4\)}
\end{problem}

\begin{answer}
A
\end{answer}

\begin{solution}
前四个命题均不恒成立. 第1个命题可取 \(b\) 为 \(x\) 轴，\(a\) 为过原点的 \(y\) 轴，\(c\) 为过原点的 \(z\) 轴，此时 \(a\perp b\)，\(b\perp c\)，但 \(a\) 与 \(c\) 不平行. 第2个命题可取 \(b\) 为 \(y\) 轴，\(a\) 为与 \(b\) 平行且位于平面 \(z=1\) 内的直线，\(c\) 为 \(x\) 轴，此时 \(b\perp c\)，但 \(a\) 与 \(c\) 异面，不能说 \(a\perp c\). 第3个命题中，取 \(\beta\) 为平面 \(z=0\)，取位于平面 \(z=1\) 内且方向与 \(x\) 轴相同的直线 \(a\)，再取平面 \(\beta\) 内方向与 \(y\) 轴相同的直线 \(b\)，则 \(a\parallel\beta\)，但 \(a\) 不平行于 \(b\). 第4个命题中，取 \(\beta\) 为平面 \(z=0\)，取平面 \(z=1\) 内的直线 \(a\) 与平面 \(z=2\) 内的直线 \(b\)，使二者方向不同，则 \(a\) 与 \(b\) 异面，且 \(b\parallel\beta\)，所以 \(b\) 不与 \(\beta\) 相交.

对于第5个命题，设 \(a\)，\(b\) 的方向向量分别为 \(\bs{u}\)，\(\bs{v}\). 因为 \(a\)，\(b\) 异面，所以 \(\bs{u}\) 与 \(\bs{v}\) 不平行，所有同时垂直于 \(a\)，\(b\) 的直线都必须与 \(\bs{u}\times\bs{v}\) 平行. 设 \(A_0\in a\)，\(B_0\in b\)，把 \(\overrightarrow{A_0B_0}\) 在由 \(\bs{u}\)，\(\bs{v}\)，\(\bs{u}\times\bs{v}\) 构成的基底下表示为 \(\overrightarrow{A_0B_0}=r\bs{u}+s\bs{v}+w(\bs{u}\times\bs{v})\). 若某条公垂线分别交 \(a\)，\(b\) 于 \(A_0+r\bs{u}\)，\(B_0-s\bs{v}\)，则其两端点之差为 \(w(\bs{u}\times\bs{v})\)，从而交点位置唯一，公垂线也唯一. 因此第5个命题为真，真命题只有 \(1\) 个，故选 A.
\end{solution}

\begin{problem}
把一同排 \(6\) 张座位编号为 \(1\)，\(2\)，\(3\)，\(4\)，\(5\)，\(6\) 的电影票全部分给 \(4\) 个人，每人至少分 \(1\) 张，至多分 \(2\) 张，且这两张票具有连续的编号，那么不同的分法种数是（\quad）

\choices
    {\(168\)}
    {\(96\)}
    {\(72\)}
    {\(144\)}
\end{problem}

\begin{answer}
D
\end{answer}

\begin{solution}
由于共有 \(6\) 张票和 \(4\) 个人，且每人至少得到 \(1\) 张、至多得到 \(2\) 张，所以恰有两个人各得到 \(2\) 张票，另外两个人各得到 \(1\) 张票. 两张票必须连续，等价于从相邻编号对
\((1,2)\)，\((2,3)\)，\((3,4)\)，\((4,5)\)，\((5,6)\)
中选出两个不共用票号的编号对. 共有 \(\mathrm C_5^2-4=6\) 种选法，其中减去的 \(4\) 种是两个编号对相邻而共用一张票的情形. 对每一种票号分组，先从四个人中选出得到两张票的两个人，再把两个连续编号对分给他们，最后把两个单票分给其余两个人，共有 \(\mathrm C_4^2\cdot2!\cdot2!=24\) 种分法. 因此总数为 \(6\cdot24=144\)，故选 D.
\end{solution}

\begin{problem}
若 \(\sin \alpha + \cos \alpha = \tan \alpha\)（\(0 < \alpha < \frac{\pi}{2}\)），则 \(\alpha \in\)（\quad）

\choices
    {\(\left(0, \frac{\pi}{6}\right)\)}
    {\(\left(\frac{\pi}{6}, \frac{\pi}{4}\right)\)}
    {\(\left(\frac{\pi}{4}, \frac{\pi}{3}\right)\)}
    {\(\left(\frac{\pi}{3}, \frac{\pi}{2}\right)\)}
\end{problem}

\begin{answer}
C
\end{answer}

\begin{solution}
令 \(g(\alpha)=\tan\alpha-\sin\alpha-\cos\alpha\). 当 \(0<\alpha\leqslant\frac\pi4\) 时，\(\tan\alpha\leqslant1<\sin\alpha+\cos\alpha\)，故方程的解不在此区间内. 又
\[
g\left(\frac\pi3\right)=\sqrt3-\frac{\sqrt3+1}{2}=\frac{\sqrt3-1}{2}>0
\]
在 \(0<\alpha<\frac\pi2\) 上，
\[
g'(\alpha)=\sec^2\alpha-\cos\alpha+\sin\alpha>1-\cos\alpha+\sin\alpha>0
\]
所以 \(g\) 严格递增，且其零点唯一，并且位于 \(\left(\frac\pi4,\frac\pi3\right)\) 内，故选 C.
\end{solution}

\begin{problem}
在函数 \( y = x^3 - 8x \) 的图象上，其切线的倾斜角小于 \( \frac{\pi}{4} \) 的点中，坐标为整数的点的个数是（\quad）

\choices
    {\(3\)}
    {\(2\)}
    {\(1\)}
    {\(0\)}
\end{problem}

\begin{answer}
D
\end{answer}

\begin{solution}
设曲线上点 \((x,y)\) 处切线的倾斜角为 \(\theta\). 由 \(y=x^3-8x\) 得切线斜率
\[
\tan\theta=y'=3x^2-8
\]
倾斜角满足 \(0\leqslant\theta<\frac\pi4\) 时，必有 \(0\leqslant3x^2-8<1\)，即
\[
\frac83\leqslant x^2<3
\]
但整数 \(x\) 的平方不可能落在区间 \([\frac83,3)\) 内，所以不存在符合条件的整点，点的个数为 \(0\)，故选 D.
\end{solution}

\begin{problem}
某初级中学有学生 \(270\) 人，其中一年级 \(108\) 人，二，三年级各 \(81\) 人. 现要利用抽样方法抽取 \(10\) 人参加某项调查，考虑选用简单随机抽样，分层抽样和系统抽样三种方案，使用简单随机抽样和分层抽样时，将学生按一，二，三年级依次统一编号为 \(1\)，\(2\)，\(\cdots\)，\(270\)；使用系统抽样时，将学生统一随机编号 \(1\)，\(2\)，\(\cdots\)，\(270\)，并将整个编号依次分为 \(10\) 段. 如果抽得号码有下列四种情况：

\begin{circlelist}
\item \(7\)，\(34\)，\(61\)，\(88\)，\(115\)，\(142\)，\(169\)，\(196\)，\(223\)，\(250\)；
\item \(5\)，\(9\)，\(100\)，\(107\)，\(111\)，\(121\)，\(180\)，\(195\)，\(200\)，\(265\)；
\item \(11\)，\(38\)，\(65\)，\(92\)，\(119\)，\(146\)，\(173\)，\(200\)，\(227\)，\(254\)；
\item \(30\)，\(57\)，\(84\)，\(111\)，\(138\)，\(165\)，\(192\)，\(219\)，\(246\)，\(270\)；
\end{circlelist}

关于上述样本的下列结论中，正确的是（\quad）

\choices
    {\circled{2}，\circled{3}都不能为系统抽样}
    {\circled{2}，\circled{4}都不能为分层抽样}
    {\circled{1}，\circled{4}都可能为系统抽样}
    {\circled{1}，\circled{3}都可能为分层抽样}
\end{problem}

\begin{answer}
D
\end{answer}

\begin{solution}
按年级编号后，三个年级对应的编号区间分别为 \([1,108]\)，\([109,189]\)，\([190,270]\). 分层抽样按人数比例抽取 \(10\) 人，应在三个年级分别抽取 \(4\)，\(3\)，\(3\) 人. 样本1和样本3的三个区间内人数均为 \(4\)，\(3\)，\(3\)，所以二者都可能是分层抽样；样本2的区间人数也是 \(4\)，\(3\)，\(3\)，而样本4的人数为 \(3\)，\(3\)，\(4\)，故样本4不可能是按比例分层抽样.

系统抽样时每段有 \(270\div10=27\) 个编号，样本必须形如 \(r\)，\(r+27\)，\(\ldots\)，\(r+9\cdot27\)，其中 \(1\leqslant r\leqslant27\). 样本1和样本3分别对应 \(r=7\) 和 \(r=11\)，均可能为系统抽样；样本2不是公差为 \(27\) 的等差数列，样本4缺少第一段 \([1,27]\) 内的编号，也不可能为系统抽样. 因此只有“1，3都可能为分层抽样”正确，故选 D.
\end{solution}

\section{填空题}
\begin{problem}
函数 \( f(x) = \frac{\sqrt{x - 2}}{x - 3} \lg \sqrt{4 - x} \) 的定义域是\fillinblank{}.
\end{problem}

\begin{answer}
\( [2,3) \cup (3,4) \)
\end{answer}

\begin{solution}
由根式有意义，需满足 \(x-2\geqslant0\)，即 \(x\geqslant2\). 因为分母不能为零，还需 \(x\ne3\). 又因为对数的真数必须为正，
\[
\sqrt{4-x}>0\Longleftrightarrow x<4
\]
综合可得定义域为 \([2,4)\)，去掉 \(3\) 后得 \([2,3)\cup(3,4)\).
\end{solution}

\begin{problem}
\(\left(x^{3} - \frac{2}{x}\right)^{4} + \left(x + \frac{1}{x}\right)^{8}\) 的展开式中整理后的常数项等于\fillinblank{}.
\end{problem}

\begin{answer}
38
\end{answer}

\begin{solution}
在 \(\left(x^3-\frac2x\right)^4\) 的展开式中，取 \(\left(-\frac2x\right)^j\) 的项，其幂次为 \(x^{3(4-j)-j}=x^{12-4j}\). 令幂次为零，得 \(j=3\)，故该部分的常数项为
\[
\mathrm C_4^3(-2)^3=-32
\]
在 \(\left(x+\frac1x\right)^8\) 的展开式中，取 \(4\) 个 \(x^{-1}\) 和 \(4\) 个 \(x\) 时幂次为零，故常数项为 \(\mathrm C_8^4=70\). 所以原展开式的常数项为 \(-32+70=38\).
\end{solution}

\begin{problem}
函数 \( y = |\sin x| \cos x - 1 \) 的最小正周期与最大值的和为\fillinblank{}.
\end{problem}

\begin{answer}
\(2\pi - \frac{1}{2}\)
\end{answer}

\begin{solution}
令 \(g(x)=|\sin x|\cos x\). 有 \(g(x+2\pi)=g(x)\)，而 \(g(x+\pi)=-g(x)\)，所以 \(\pi\) 不是它的周期. 函数的零点为 \(k\frac\pi2\)（\(k\in\mathbb Z\)），若存在小于 \(2\pi\) 的正周期，则该周期必须是 \(\frac\pi2\)，\(\pi\) 或 \(\frac{3\pi}{2}\). 但 \(g\left(\frac\pi4\right)=\frac12\ne g\left(\frac{3\pi}{4}\right)=-\frac12\)，故 \(\frac\pi2\) 不是周期；又 \(g\left(\frac\pi4\right)=\frac12\ne g\left(\frac{5\pi}{4}\right)=-\frac12\)，故 \(\pi\) 不是周期；最后 \(g\left(\frac{3\pi}{4}\right)=-\frac12\ne g\left(\frac{9\pi}{4}\right)=\frac12\)，故 \(\frac{3\pi}{2}\) 也不是周期. 所以最小正周期为 \(2\pi\).

又因为
\[
g(x)=|\sin x|\cos x\leqslant|\sin x\cos x|\leqslant\frac12
\]
且在 \(x=\frac\pi4\) 时取到等号，所以原函数的最大值为 \(\frac12-1=-\frac12\). 所求和为 \(2\pi-\frac12\).
\end{solution}

\begin{problem}
某实验室需购某种化工原料 \(106\) 千克，现在市场上该原料有两种包装，一种是每袋 \(35\) 千克，价格为 \(140\) 元；另一种是每袋 \(24\) 千克，价格为 \(120\) 元. 在满足需要的条件下，最少要花费\fillinblank{} 元.
\end{problem}

\begin{answer}
500
\end{answer}

\begin{solution}
设购买每袋 \(35\) 千克和每袋 \(24\) 千克的原料分别为 \(x\)，\(y\) 袋，其中 \(x\)，\(y\) 为非负整数. 需满足
\[
35x+24y\geqslant106
\]
费用为 \(140x+120y\) 元. 先取 \(x=1\)，\(y=3\)，原料质量为 \(35+3\times24=107\) 千克，费用为 \(140+3\times120=500\) 元.

下面说明费用不能低于 \(500\) 元. 若费用小于 \(500\) 元，则 \(x\leqslant3\). 当 \(x=0\) 时，至多取 \(y=4\)，质量不超过 \(96\) 千克；当 \(x=1\) 时，至多取 \(y=2\)，质量不超过 \(83\) 千克；当 \(x=2\) 时，至多取 \(y=1\)，质量不超过 \(94\) 千克；当 \(x=3\) 时，至多取 \(y=0\)，质量为 \(105\) 千克. 这些情况都不能满足 \(106\) 千克的需要，故最少花费为 \(500\) 元.
\end{solution}

\section{解答题}
\begin{problem}
已知向量 \(\bs{a} = (x^2, x + 1)\)，\(\bs{b} = (1 - x, t)\)，若函数 \(f(x) = \bs{a} \cdot \bs{b}\) 在区间 \((-1, 1)\) 上是增函数，求 \(t\) 的取值范围.
\end{problem}

\begin{solution}
由向量数量积的坐标运算，得
\[
f(x)=x^2(1-x)+t(x+1)=-x^3+x^2+tx+t
\]
因此
\[
f'(x)=-3x^2+2x+t
\]
若 \(f(x)\) 在 \((-1,1)\) 上为增函数，则在该区间内必有 \(f'(x)\geqslant0\). 当 \(t<5\) 时，\(f'(-1)=t-5<0\)，由导函数的连续性可知，在 \(-1\) 的右侧存在一段区间使 \(f'(x)<0\)，这与增函数矛盾，所以 \(t\geqslant5\).

当 \(t\geqslant5\) 时，对任意 \(x\in(-1,1)\)，有
\[
f'(x)\geqslant-3x^2+2x+5=-3(x+1)\left(x-\frac{5}{3}\right)>0
\]
所以 \(f(x)\) 在 \((-1,1)\) 上为增函数. 故 \(t\) 的取值范围是 \([5,+\infty)\).
\end{solution}

\begin{problem}
在 \(\triangle ABC\) 中，已知 \(\tan B = \sqrt{3}\)，\(\cos C = \frac{1}{3}\)，\(AC = 3\sqrt{6}\)，求 \(\triangle ABC\) 的面积.
\end{problem}

\begin{solution}
设 \(a=BC\)，\(b=CA\)，\(c=AB\). 由于 \(B\) 是三角形内角且 \(\tan B=\sqrt{3}\)，所以 \(B=\dfrac{\pi}{3}\). 又 \(\cos C=\dfrac{1}{3}\) 且 \(C\) 为三角形内角，故
\[
\sin C=\sqrt{1-\cos^2C}=\frac{2\sqrt{2}}{3}
\]
由正弦定理，
\[
\frac{b}{\sin B}=\frac{3\sqrt{6}}{\sqrt{3}/2}=6\sqrt{2}
\]
从而
\[
c=6\sqrt{2}\sin C=6\sqrt{2}\cdot\frac{2\sqrt{2}}{3}=8
\]
三角形内角和给出 \(A=\pi-B-C\)，因此
\[
\sin A=\sin(B+C)=\sin B\cos C+\cos B\sin C
 =\frac{\sqrt{3}}{6}+\frac{\sqrt{2}}{3}
 =\frac{\sqrt{3}+2\sqrt{2}}{6}
\]
于是
\[
S_{\triangle ABC}=\frac{1}{2}bc\sin A
 =\frac{1}{2}\cdot3\sqrt{6}\cdot8\cdot\frac{\sqrt{3}+2\sqrt{2}}{6}
 =6\sqrt{2}+8\sqrt{3}
\]
\end{solution}

\begin{problem}
设数列 \(\{a_{n}\}\) 的前 \(n\) 项和为 \(S_{n} = 2n^{2}\)，\(\{b_{n}\}\) 为等比数列，且 \(a_{1} = b_{1}\)，\(b_{2}(a_{2} - a_{1}) = b_{1}\).

\begin{enumerate}
\item 求数列 \(\{a_{n}\}\) 和 \(\{b_{n}\}\) 的通项公式；
\item 设 \(c_{n} = \frac{a_{n}}{b_{n}}\)，求数列 \(\{c_{n}\}\) 的前 \(n\) 项和 \(T_{n}\).
\end{enumerate}
\end{problem}

\begin{solution}
\begin{enumerate}
\item 由 \(S_1=a_1\)，得 \(a_1=2\). 当 \(n\geqslant2\) 时，
\[
a_n=S_n-S_{n-1}=2n^2-2(n-1)^2=4n-2
\]
当 \(n=1\) 时，\(4n-2=2=a_1\)，所以数列 \(\{a_n\}\) 的通项公式为 \(a_n=4n-2\). 由 \(a_1=b_1\)，得 \(b_1=2\)，又 \(a_2=6\)，代入条件得 \(b_2(6-2)=2\)，从而 \(b_2=\frac12\). 因此等比数列 \(\{b_n\}\) 的公比为
\[
q=\frac{b_2}{b_1}=\frac14
\]
所以 \(b_n=2\left(\frac14\right)^{n-1}\).
\item 由上面的通项公式，
\[
c_n=\frac{a_n}{b_n}=\frac{4n-2}{2(1/4)^{n-1}}=(2n-1)4^{n-1}
\]
令 \(T_n=\sum_{k=1}^{n}c_k\). 注意到
\[
(6k-5)4^k-(6k-11)4^{k-1}=9(2k-1)4^{k-1}
\]
对 \(k=1\)，\(2\)，\(\ldots\)，\(n\) 求和，左端相邻两项依次相消，得
\[
9T_n=(6n-5)4^n-(-5)=(6n-5)4^n+5
\]
故 \(T_n=\frac{(6n-5)4^n+5}{9}\).
\end{enumerate}
\end{solution}

\begin{problem}
如图所示的多面体是由底面为 \(ABCD\) 的长方体被截面 \(AEC_{1}F\) 所截得的，其中 \(AB = 4\)，\(BC = 2\)，\(CC_{1} = 3\)，\(BE = 1\).

\begin{enumerate}
\item 求 \(BF\) 的长；
\item 求点 \(C\) 到平面 \(AEC_{1}F\) 的距离.
\end{enumerate}

\bitmapfigure[max width=0.56\linewidth]{img/2005/hubei_liberal/q20_fig1.png}
\end{problem}

\begin{solution}
\begin{enumerate}
\item 以 \(D\) 为原点，分别以 \(DA\)，\(DC\)，\(DD_1\) 所在直线为 \(x\) 轴，\(y\) 轴，\(z\) 轴建立空间直角坐标系，则
\(D(0,0,0)\)，\(A(2,0,0)\)，\(B(2,4,0)\)，\(C(0,4,0)\)，\(C_1(0,4,3)\).
由于 \(E\) 在 \(BB_1\) 上且 \(BE=1\)，所以 \(E(2,4,1)\). 设截面平面为 \(\alpha\). 平面 \(ABB_1A_1\) 与平面 \(DCC_1D_1\) 平行，且它们与 \(\alpha\) 的交线分别为 \(AE\) 与 \(FC_1\)，所以 \(AE\parallel FC_1\). 平面 \(ADD_1A_1\) 与平面 \(BCC_1B_1\) 平行，且它们与 \(\alpha\) 的交线分别为 \(AF\) 与 \(EC_1\)，所以 \(AF\parallel EC_1\). 因此四边形 \(AEC_1F\) 是平行四边形. 设 \(F(0,0,z)\)，则
\(\overrightarrow{AF}=(-2,0,z)\)，\(\overrightarrow{EC_1}=(-2,0,2)\).
由 \(\overrightarrow{AF}=\overrightarrow{EC_1}\) 得 \(z=2\)，即 \(F(0,0,2)\).
由点的坐标，
\[
BF=\sqrt{(2-0)^2+(4-0)^2+(0-2)^2}=\sqrt{24}=2\sqrt6
\]
\item 由 \(\overrightarrow{AE}=(0,4,1)\)，\(\overrightarrow{AF}=(-2,0,2)\)，可取平面 \(AEC_1F\) 的一个法向量为
\[
\bs{n}=\overrightarrow{AE}\times\overrightarrow{AF}=(8,-2,8)\sim(4,-1,4)
\]
所以平面 \(AEC_1F\) 的方程为
\[
4(x-2)-y+4z=0
\]
即 \(4x-y+4z-8=0\). 因此点 \(C(0,4,0)\) 到该平面的距离为
\[
d=\frac{|4\cdot0-4+4\cdot0-8|}{\sqrt{4^2+(-1)^2+4^2}}=\frac{12}{\sqrt{33}}=\frac{4\sqrt{33}}{11}
\]
\end{enumerate}
\end{solution}

\begin{problem}
某会议室用 \(5\) 盏灯照明，每盏灯各使用灯泡一只，且型号相同. 假定每盏灯能否正常照明只与灯泡的寿命有关，该型号的灯泡寿命为 \(1\) 年以上的概率为 \(p_1\)，寿命为 \(2\) 年以上的概率为 \(p_2\). 从使用之日起每满 \(1\) 年进行一次灯泡更换工作，只更换已坏的灯泡，平时不换.

\begin{enumerate}
\item 在第一次灯泡更换工作中，求不需要更换灯泡的概率和更换 \(2\) 只灯泡的概率；
\item 在第二次灯泡更换工作中，对其中的某一盏灯来说，求该盏灯需要更换灯泡的概率；
\item 当 \(p_1 = 0.8\)，\(p_2 = 0.3\) 时，求在第二次灯泡更换工作中，至少需要更换 \(4\) 只灯泡的概率（结果保留两个有效数字）.
\end{enumerate}
\end{problem}

\begin{solution}
\begin{enumerate}
\item 假设各灯泡的寿命相互独立，且更换后的灯泡与初始灯泡同型号. 第一次更换工作时，某只初始灯泡不需要更换的概率为 \(p_1\). 五只灯泡都不需要更换，必须五只都能使用 \(1\) 年以上，因此概率为 \(p_1^5\). 恰好更换 \(2\) 只灯泡时，从五只灯泡中选出这 \(2\) 只，且其余 \(3\) 只均能使用 \(1\) 年以上，所以概率为
\[
\mathrm C_5^2(1-p_1)^2p_1^3
\]
\item 考虑其中一盏灯. 若初始灯泡使用不到 \(1\) 年就坏掉，第一次更换后装上的新灯泡又使用不到 \(1\) 年，则第二次更换时需要更换灯泡，其概率为 \((1-p_1)^2\). 若初始灯泡能使用 \(1\) 年以上但不能使用 \(2\) 年，则第二次更换时也需要更换灯泡，其概率为 \(p_1-p_2\). 这两种情况互斥，所以第二次更换时该灯需要更换灯泡的概率为
\[
(1-p_1)^2+p_1-p_2
\]
\item 当 \(p_1=0.8\)，\(p_2=0.3\) 时，第二次更换时一盏灯需要更换灯泡的概率为
\[
p=(1-0.8)^2+0.8-0.3=0.54
\]
五盏灯中需要更换的灯数服从二项分布. 至少需要更换 \(4\) 只灯泡包括恰好更换 \(4\) 只和恰好更换 \(5\) 只两种情况，故所求概率为
\[
\mathrm C_5^4p^4(1-p)+p^5=\mathrm C_5^4(0.54)^4(0.46)+(0.54)^5=0.2414867904\approx0.24
\]
\end{enumerate}
\end{solution}

\begin{problem}
设 \(A\)、\(B\) 是椭圆 \(3x^{2} + y^{2} = \lambda\) 上的两点，点 \(N(1,3)\) 是线段 \(AB\) 的中点，线段 \(AB\) 的垂直平分线与椭圆相交于 \(C\)、\(D\) 两点.

\begin{enumerate}
\item 确定 \(\lambda\) 的取值范围，并求直线 \(AB\) 的方程；
\item 试判断是否存在这样的 \(\lambda\)，使得 \(A\)、\(B\)、\(C\)、\(D\) 四点在同一个圆上？并说明理由.
\end{enumerate}
\end{problem}

\begin{solution}
\begin{enumerate}
\item 直线 \(AB\) 不可能是竖直线，因为若 \(x=1\)，椭圆上的两个交点的纵坐标之和为 \(0\)，其中点纵坐标为 \(0\)，不可能是 \(N(1,3)\). 设直线 \(AB\) 的斜率为 \(k\)，则其方程为
\[
y-3=k(x-1)
\]
将 \(y=kx+3-k\) 代入椭圆方程，得关于 \(x\) 的方程
\[
(3+k^2)x^2+2k(3-k)x+(3-k)^2-\lambda=0
\]
设交点 \(A\)，\(B\) 的横坐标分别为 \(x_1\)，\(x_2\). 因为 \(N(1,3)\) 是线段 \(AB\) 的中点，所以 \(x_1+x_2=2\). 由韦达定理，
\[
-\frac{2k(3-k)}{3+k^2}=2
\]
即 \(-2k(3-k)=2(3+k^2)\)，从而 \(k=-1\). 所以直线 \(AB\) 的方程为
\[
x+y-4=0
\]
令 \(A(1+s,3-s)\)，\(B(1-s,3+s)\)，将这两点代入椭圆方程，得
\[
3(1+s)^2+(3-s)^2=12+4s^2=\lambda
\]
因为 \(A\)，\(B\) 是两个不同的交点，故 \(s\ne0\)，从而 \(\lambda>12\). 反之，当 \(\lambda>12\) 时，\(s=\frac{\sqrt{\lambda-12}}2\) 为非零实数，确有两个这样的交点. 因此 \(\lambda\) 的取值范围为 \(\lambda\in(12,+\infty)\).
\item 直线 \(AB\) 的斜率为 \(-1\)，所以其垂直平分线 \(CD\) 的斜率为 \(1\)，且过 \(N(1,3)\)，故
\[
y-3=x-1
\]
即 \(y=x+2\). 设 \(C(1+t,3+t)\)，\(D(1+t',3+t')\)，则 \(t\)，\(t'\) 是方程
\[
4t^2+12t+12-\lambda=0
\]
的两个根. 当 \(\lambda>12\) 时，该方程的判别式为 \(16(\lambda-3)>0\)，所以 \(C\)，\(D\) 确为两个不同的交点.

考虑方程
\[
(x-1)^2+(y-3)^2+3(x+y-4)-\frac{\lambda-12}{2}=0
\]
它等价于
\[
\left(x+\frac{1}{2}\right)^2+\left(y-\frac{3}{2}\right)^2=\frac{\lambda-3}{2}
\]
当 \(\lambda>12\) 时，右端为正，故该方程表示一个圆. 对点 \(A\)，\(B\)，有 \((x-1)^2+(y-3)^2=2s^2=\frac{\lambda-12}{2}\) 且 \(x+y-4=0\)，所以 \(A\)，\(B\) 在该圆上.

对点 \(C\) 或 \(D\)，令对应的参数为 \(t\) 或 \(t'\). 由 \(4t^2+12t+12-\lambda=0\)，得
\[
2t^2+6t-\frac{\lambda-12}{2}=0
\]
又 \((x-1)^2+(y-3)^2=2t^2\)，\(x+y-4=2t\)，所以 \(C\)，\(D\) 也都满足所考虑的圆方程. 故存在这样的 \(\lambda\)，事实上对任意 \(\lambda>12\)，\(A\)，\(B\)，\(C\)，\(D\) 四点均在同一个圆上.
\end{enumerate}
\end{solution}
